\qhead{5. Choosing Prices Levels — Which price levels does your company apply? What are the main considerations in that decision?}

BYD's distinctive choice is to run \textbf{different price-level strategies in different segments at once} \parencite{nagle2018}.

\textbf{Penetration in the mass market (Seagull, Dynasty).} Every condition holds: highly \emph{elastic} demand, large \emph{scale and learning} economies from vertical integration, and explicit \emph{market-share} goals. Crucially, pricing near cost is rational here \emph{because BYD's cost is genuinely below rivals'} — low price does not mean below cost.

\textbf{Skimming at the top (Yangwang).} Near-exchange-value pricing — the U9 at about \$3\,million \parencite{electrekU9}. Demand is inelastic and the purpose is to showcase technology and recover R\&D.

\textbf{Neutral-to-skim in Europe.} BYD prices just under Tesla and VW, capturing a tariff-protected margin (Question~12).

\textbf{Defining the window.} The floor is incremental cost (lowered by the Blade Battery); the ceiling is exchange value versus the local reference car. The structural cost edge \textbf{widens the window}, giving room to choose penetration without going below cost.

\textbf{Main consideration / risk.} Penetration ``trains customers to expect low prices'' and can become unsustainable: 2025 net profit fell 19\% despite record volume, and Chinese regulators intervened \parencite{byd2025}. The skim end yields prestige but no scale.
